\subsection{\texorpdfstring{第~\ref{sec:lie-groups}~节}{Lie 群}}\label{subsec:solutions-lie-groups}

\begin{description}[font=\normalfont,leftmargin=2em,labelwidth=1.5em]
\RGsolution{ex:affine-maps-lie-group}{affine-maps-lie-group}
\begin{enumerate}[label=(\alph*),ref=\alph*,leftmargin=*]
\item 给定两个仿射映射 $g(t)=yt+x$ 和 $h(t)=wt+z$, 有
\[
  (g\circ h)(t)=g(h(t))=g(wt+z)=ywt+yz+x.
\]
因此, 群运算由
\[
  (x,y)\cdot(z,w)=(yz+x,yw)
\]
给出. 单位元显然是 $e=(0,1)$ (对应于恒等映射), 故
\begin{align*}
  (z,w)=(x,y)^{-1}&\iff(yz+x,yw)=(0,1)\\
  &\iff(z,w)=\left(-\frac{x}{y},\frac1y\right).
\end{align*}
因此映射 $H\times H\ni(g,h)\mapsto g\cdot h\in H$ 与 $H\ni g\mapsto g^{-1}\in H$ 都是光滑的, 所以 $H$ 是 Lie 群.
\item 因为
\[
  L_{(x,y)}(z,w)=(yz+x,yw),
\]
所以 $(\mathrm{d}L_{(x,y)})_{(z,w)}$ 的矩阵表示为
\[
  (\mathrm{d}L_{(x,y)})_{(z,w)}=\begin{pmatrix}y&0\\0&y\end{pmatrix}.
\]
因此 $X^V_{(x,y)}$ 的分量为
\[
  \begin{pmatrix}y&0\\0&y\end{pmatrix}\begin{pmatrix}\xi\\\eta\end{pmatrix}
  =\begin{pmatrix}y\xi\\y\eta\end{pmatrix}.
\]
\item\label{sol:affine-group-bracket} 若
\[
  V=\xi\frac{\partial}{\partial x}+\eta\frac{\partial}{\partial y}
  \quad\text{且}\quad
  W=\zeta\frac{\partial}{\partial x}+\omega\frac{\partial}{\partial y},
\]
则
\begin{align*}
  [X^V,X^W]
  &=\left[\xi y\frac{\partial}{\partial x}+\eta y\frac{\partial}{\partial y},
    \zeta y\frac{\partial}{\partial x}+\omega y\frac{\partial}{\partial y}\right]\\
  &=(\eta\zeta-\omega\xi)y\frac{\partial}{\partial x}.
\end{align*}
因此
\[
  [V,W]=[X^V,X^W]_{(0,1)}=(\eta\zeta-\omega\xi)\frac{\partial}{\partial x}.
\]
\item\label{sol:affine-group-exponential} $X^V$ 的流由常微分方程组
\[
  \begin{cases}
  \dot x=\xi y,\\
  \dot y=\eta y
  \end{cases}
\]
的解给出. 当 $\eta\ne0$ 时, 解为
\[
  \begin{cases}
  x=x_0+y_0\dfrac{\xi(e^{\eta t}-1)}{\eta},\\[2pt]
  y=y_0e^{\eta t},
  \end{cases}
\]
而当 $\eta=0$ 时, 解为
\[
  \begin{cases}
  x=x_0+y_0\xi t,\\
  y=y_0.
  \end{cases}
\]
令 $(x_0,y_0)=e=(0,1)$ 且 $t=1$, 就得到指数映射: 当 $\eta\ne0$ 时,
\[
  \exp(V)=\left(\frac{\xi(e^\eta-1)}{\eta},e^\eta\right),
\]
而当 $\eta=0$ 时,
\[
  \exp(V)=(\xi,1).
\]
\item 两个这样的矩阵相乘为
\[
  \begin{pmatrix}y&x\\0&1\end{pmatrix}\begin{pmatrix}w&z\\0&1\end{pmatrix}
  =\begin{pmatrix}yw&yz+x\\0&1\end{pmatrix},
\]
这恰好给出 $H$ 上的群运算. 因此, 可以把 $H$ 与 $GL(2)$ 的相应子群等同. 满足 $c(0)=I$ 的曲线 $c:(-\varepsilon,\varepsilon)\to H$ 此时由
\[
  c(t)=\begin{pmatrix}y(t)&x(t)\\0&1\end{pmatrix},
\]
给出, 其中 $x(0)=0$ 且 $y(0)=1$, 它在 $t=0$ 时的导数是
\[
  \dot c(0)=\begin{pmatrix}\dot y(0)&\dot x(0)\\0&0\end{pmatrix}.
\]
因此, $\mathfrak{h}$ 可以与形如
\[
  \begin{pmatrix}\eta&\xi\\0&0\end{pmatrix}
\]
的矩阵所构成的向量空间等同. Lie 括号于是必由
\begin{align*}
  \left[\begin{pmatrix}\eta&\xi\\0&0\end{pmatrix},\begin{pmatrix}\omega&\zeta\\0&0\end{pmatrix}\right]
  &=\begin{pmatrix}\eta&\xi\\0&0\end{pmatrix}\begin{pmatrix}\omega&\zeta\\0&0\end{pmatrix}
    -\begin{pmatrix}\omega&\zeta\\0&0\end{pmatrix}\begin{pmatrix}\eta&\xi\\0&0\end{pmatrix}\\
  &=\begin{pmatrix}0&\eta\zeta-\omega\xi\\0&0\end{pmatrix},
\end{align*}
给出, 这与 (\ref{sol:affine-group-bracket}) 一致. 此外, 指数映射必由
\begin{align*}
  \exp\begin{pmatrix}\eta&\xi\\0&0\end{pmatrix}
  &=\sum_{k=0}^{+\infty}\frac{1}{k!}\begin{pmatrix}\eta&\xi\\0&0\end{pmatrix}^k\\
  &=\begin{pmatrix}1&0\\0&1\end{pmatrix}+\begin{pmatrix}\eta&\xi\\0&0\end{pmatrix}
      +\frac12\begin{pmatrix}\eta^2&\eta\xi\\0&0\end{pmatrix}+\cdots,
\end{align*}
给出, 从而当 $\eta\ne0$ 时,
\[
  \exp\begin{pmatrix}\eta&\xi\\0&0\end{pmatrix}
  =\begin{pmatrix}e^\eta&\dfrac{\xi(e^\eta-1)}{\eta}\\0&1\end{pmatrix},
\]
而当 $\eta=0$ 时,
\[
  \exp\begin{pmatrix}\eta&\xi\\0&0\end{pmatrix}=\begin{pmatrix}1&\xi\\0&1\end{pmatrix},
\]
这与 (\ref{sol:affine-group-exponential}) 一致.
\end{enumerate}

\RGsolution{ex:determinant-of-matrix-exponential}{determinant-of-matrix-exponential}
\begin{enumerate}[label=(\alph*),ref=\alph*,leftmargin=*]
\item 显然 $h(\mathbb{R})=\mathbb{R}\setminus\{0\}$, 且
\begin{align*}
  h(x+y)&=\det e^{A(x+y)}=\det(e^{Ax}e^{Ay})\\
        &=(\det e^{Ax})(\det e^{Ay})=h(x)\cdot h(y).
\end{align*}
\item 注意 $h(t)=f(e^{At})$, 其中 $f:GL(n)\to\mathbb{R}$ 由 $f(B)=\det(B)$ 给出. 因此, 由于 $(e^{At})'=Ae^{At}$, 由 Example~\ref{ex:classical-matrix-lie-groups}(\ref{ex:special-linear-group-regular-level-set}) 得
\[
  h'(0)=(\mathrm{d}f)_I(A)=\operatorname{tr}(A).
\]
\item 再由 Example~\ref{ex:classical-matrix-lie-groups}(\ref{ex:special-linear-group-regular-level-set}), 有
\[
  h'(t)
  =(\mathrm{d}f)_{e^{At}}(Ae^{At})
  =\det(e^{At})\operatorname{tr}(e^{-At}Ae^{At})
  =h(t)\operatorname{tr}(A).
\]
所以 $h(t)=ke^{\operatorname{tr}(A)t}$, 其中 $k\in\mathbb{R}$ 是常数. 已知 $h(0)=\det I=1$, 故 $h(t)=e^{\operatorname{tr}(A)t}$, 从而 $\det(e^A)=h(1)=e^{\operatorname{tr}(A)}$.
\end{enumerate}

\RGsolution{ex:differential-of-lie-group-homomorphism}{differential-of-lie-group-homomorphism}
\begin{enumerate}[label=(\alph*),ref=\alph*,leftmargin=*]
\item 设 $V\in\mathfrak{g}$. 先证明 $X^V$ 与 $X^{(\mathrm{d}F)_eV}$ 关于 $F$ 相关 [参见 Exercise~\ref{ex:vector-fields}(\ref{ex:related-vector-fields})]. 的确,
\begin{align*}
  (\mathrm{d}F)_gX_g^V
  &=(\mathrm{d}F)_g(\mathrm{d}L_g)_eV=\mathrm{d}(F\circ L_g)_eV\\
  &=\mathrm{d}(L_{F(g)}\circ F)_eV=(\mathrm{d}L_{F(g)})_e(\mathrm{d}F)_eV\\
  &=X_{F(g)}^{(\mathrm{d}F)_eV},
\end{align*}
其中 $F\circ L_g=L_{F(g)}\circ F$ 是因为 $F$ 为 Lie 群同态.

设 $V,W\in\mathfrak{g}$. 则
\begin{align*}
  (\mathrm{d}F)_e[V,W]&=(\mathrm{d}F)_e[X^V,X^W]_e=(F_*[X^V,X^W])_e\\
  &=[F_*X^V,F_*X^W]_e=[X^{(\mathrm{d}F)_eV},X^{(\mathrm{d}F)_eW}]_e\\
  &=[(\mathrm{d}F)_eV,(\mathrm{d}F)_eW].
\end{align*}
这里使用了 Exercise~\ref{ex:vector-fields}(\ref{ex:pushforward-preserves-lie-bracket}) 的结论, 它很容易推广到任意可微映射.
\item 给定 $g\in G$, 有 $F=L_{F(g)}\circ F\circ L_{g^{-1}}$. 因此,
\[
  (\mathrm{d}F)_g=(\mathrm{d}L_{F(g)})_e\circ(\mathrm{d}F)_e\circ(\mathrm{d}L_{g^{-1}})_g.
\]
由于左乘映射是微分同胚, 若 $(\mathrm{d}F)_e$ 是同构, 则对任意 $g\in G$, $(\mathrm{d}F)_g$ 都是同构. 由反函数定理, $F$ 是局部微分同胚.
\item 设 $U\ni e$ 是开集, 使 $F$ 限制在 $U$ 上是到 $F(U)$ 的微分同胚. 对任意 $g\in G$, $F=L_{F(g)}\circ F\circ L_{g^{-1}}$ 限制在 $L_g(U)$ 上是到其像的微分同胚. 此外, 容易看出 $F^{-1}(F(L_g(U)))=L_g(F^{-1}(F(U)))$. 因此, 只需检验 $F^{-1}(F(U))$ 是若干与 $F(U)$ 微分同胚的开集的不交并. $F(h)\in F(U)$ 当且仅当对某个 $h_0\in U$ 有 $F(h)=F(h_0)$, 即当且仅当 $hh_0^{-1}=g\in\ker(F)$. 因此
\[
  F^{-1}(F(U))=\bigcup_{g\in\ker(F)}L_g(U).
\]
另一方面, 开集 $L_g(U)$ 显然互不相交, 因为若 $g\in\ker(F)$ 且 $h_0\in U$, 则 $F(gh_0)=F(h_0)$. 因此, 由于 $F$ 在 $U$ 上是单射, 只有当 $g=e$ 时 $gh_0$ 才能在 $U$ 中.
\end{enumerate}
\end{description}
